\section{Conclusion}
\label{sec:conclusion}

\subsection{Summary}

We have presented a candidate algorithmic toolbox and requirements map for
congressional-redistricting research. The implementation and evidence described
here do not by themselves satisfy the eleven legal requirements or establish a
post-\callais{} compliance framework.

The toolbox is organised along three orthogonal axes:
Axis A controls what METIS balances (population only, population plus area,
population with VRA alignment, population with partisan vote share);
Axis B controls edge management (geographic boundary lengths, county-sticky
weights, unweighted);
Axis C controls division strategy (balanced bisection, GeoSection, AreaSection,
ApportionRegions, NestSection).

Each legal requirement maps to at most one component on one axis.
The components compose freely except for the Callais mutex: VRASection (Axis~A3)
and ProportionalSection (Axis~A4) are mutually exclusive, enforced at the CLI
level by \texttt{callais\_preflight} and documented in the plan manifest.

The four-State bakeoff retains confirmed, estimated, and pending cells. It
therefore generates hypotheses about geographic sorting, algorithm choice,
minority-opportunity diagnostics, subdivision tradeoffs, and audit structure;
it does not establish cross-State effects, causal mechanisms, legal compliance,
or algorithm rankings.
The separately preregistered Wisconsin proof slice establishes a narrower
implementation history. Its initial two validity failures exposed hidden
tolerance, reporting, and contiguity-path defects. After remediation, all four
plans passed their native audits and regenerated exactly. The subsequently
frozen eight-State national pilot retained a second failure witness, exposed an
equal-split recursive contiguity defect, and passed and exactly regenerated all
32 cells after remediation. The pilot cannot
support national prevalence, causal effects, or a ranking, and the frozen
44-State clean matrix subsequently passed all 176 cells without seed retries,
and an independent 176-cell execution regenerated the normalized evidence
exactly. Its scope is
implementation validity rather than comparative superiority or legal adequacy.

\subsection{The Callais Compliance Design}

\callais{}'s disentanglement requirement at p.~36 is the structurally decisive
constraint for algorithmic redistricting.
The proposed toolbox is designed to separate signal types: each mode uses exactly one signal type,
the two signal-bearing modes are mutex-gated, and every run is documented in
a tamper-detectable plan manifest.

The strong-inference framework of \callais{} holding~2 is directly supported:
a practitioner can run GeoSection (pure political goals) and VRASection (minority
alignment) independently, compare the results, and demonstrate that the political
goals could have been achieved with different threshold-based minority outcomes.
For Alabama, this is an illustrative research comparison, not confirmation of
an \allenmilligan{} remedy.

\subsection{The Central Empirical Finding}

The central hypothesis is that geographic voter sorting materially contributes
to partisan redistricting outcomes. This matrix does not estimate that effect
separately from algorithm, profile, resolution, seed, and repair choices, and it
does not support a claim that outcomes are invariant to district lines.

This finding has a constructive implication.
The evidence here cannot adjudicate a claim that an algorithm favoured a party.
A relevant empirical question is whether another valid, comparably constrained
method could achieve the same stated goals with different minority outcomes,
but answering it requires a completed counterfactual design rather than this
mixed-evidence matrix or the one-State mechanics slice.
For \allenmilligan{} matters, VRASection may help construct an illustrative
diagnostic; it does not provide the legal or factual proof required by a court.

\subsection{Future Work}

Four extensions are needed to complete the toolbox:

\begin{enumerate}
  \item \textbf{T.4 ApportionRegions}: Extend the completed, exactly regenerated
    frozen-cohort validity result only under separately preregistered comparative
    designs, preserving every failure and separating validity from performance.
  \item \textbf{T.5 ProportionalSection}: Implementation with the
    Callais-compliant mutex gate; empirical comparison against AreaSection
    in state-court proportionality challenge states.
  \item \textbf{T.3 Phase~2--3}: Download and spatial join of TIGER MCD, place,
    and VTD files; parameter sweep at the township and precinct levels for
    MN, WI, and New England states.
  \item \textbf{T.8 Three-Census CSS}: Completion of the 2010 GeoSection sweep
    to compute full three-census CSS for all 50 states.
\end{enumerate}

These extensions are prerequisites for a broader auditable research toolbox;
they would not by themselves establish a legally sufficient framework for a
future redistricting cycle.

\paragraph{Data Availability.}
The \texttt{bisect} Rust binary implementing GeoSection, VRASection,
AreaSection, NestSection, and StabilitySection is available as open-source
software at \url{https://github.com/apportionment-research/redist}.
All redistricting outputs, adjacency graphs, and analysis scripts are available
at the same repository.
Pre-computed adjacency files for all 50 states are available as GitHub Release assets,
with SHA-256 hashes recorded in the plan manifests.
All parameters are documented in \texttt{docs/REDIST\_CLI.md}.
Marked empirical results are intended to be reproducible from their stated
parameters and inputs. Estimated and pending cells are not execution results.
